\SourceLeaf{89}
\Article{1}
Zur Erklärung der magnetischen Erscheinungen setzen wir zwei magnetische Flüssigkeiten voraus. Mit den Physikern nennen wir die eine die nördliche, die andere die südliche. Wir nehmen an, dass die Elemente der einen Flüssigkeit die der anderen anziehen, dass sich dagegen je zwei Elemente derselben Flüssigkeit gegenseitig abstoßen und dass beide Wirkungen umgekehrt proportional zum Quadrat der Entfernung sind. Weiter unten wird sich zeigen, dass unsere eigenen Beobachtungen die Richtigkeit dieses Gesetzes vollständig bestätigen.

Diese Flüssigkeiten treten für sich allein nicht in Erscheinung, sondern nur in Verbindung mit den wägbaren Teilchen von Körpern, die Magnetismus annehmen können. Ihre Wirkungen äußern sich darin, dass sie diese Teilchen zur Bewegung antreiben oder die Bewegung verhindern oder verändern, die andere auf die Teilchen wirkende Kräfte, etwa die Schwerkraft, hervorrufen würden.

Die Wirkung einer gegebenen Menge magnetischer Flüssigkeit auf eine gegebene Menge derselben oder der anderen Flüssigkeit in einer gegebenen Entfernung lässt sich demnach mit einer gegebenen bewegenden Kraft vergleichen, also mit der Wirkung einer gegebenen beschleunigenden Kraft auf eine gegebene Masse. Da wir die magnetischen Flüssigkeiten selbst nur an den Wirkungen erkennen können, die sie hervorrufen, müssen eben diese Wirkungen zu ihrer Messung dienen.

Um dieses Maß auf eindeutige Begriffe zurückzuführen, müssen wir zunächst für drei Arten von Größen Einheiten festlegen, nämlich für Entfernungen, wägbare Massen und beschleunigende Kräfte. Als dritte Einheit kann man die Schwerkraft am Beobachtungsort wählen. Erscheint diese Wahl weniger geeignet, so muss noch eine Zeiteinheit hinzukommen. Dann ist für uns diejenige beschleunigende Kraft $=1$, die in einer Zeiteinheit die Geschwindigkeit eines Körpers, der sich in ihrer Richtung bewegt, um die Einheit der Geschwindigkeit ändert.

\SourceLeaf{90}
Unter diesen Festlegungen ist die Einheitsmenge der nördlichen Flüssigkeit diejenige Menge, deren abstoßende Wirkung auf eine gleich große Menge in der Entfernung $=1$ einer bewegenden Kraft $=1$ entspricht, also der Wirkung einer beschleunigenden Kraft $=1$ auf eine Masse $=1$. Dasselbe gilt für die Einheitsmenge der südlichen Flüssigkeit. Bei dieser Bestimmung muss man sich offenbar sowohl die wirkende Flüssigkeit als auch die Flüssigkeit, auf die sie wirkt, in physischen Punkten konzentriert denken. Außerdem ist anzunehmen, dass die Anziehung zwischen gegebenen Mengen ungleichnamiger Flüssigkeiten in einer gegebenen Entfernung der Abstoßung zwischen jeweils gleich großen Mengen gleichnamiger Flüssigkeiten entspricht. Die Wirkung einer Menge $m$ nördlicher magnetischer Flüssigkeit auf eine Menge $m'$ derselben Flüssigkeit in der Entfernung $r$ wird also, wenn man beide jeweils in einem physischen Punkt konzentriert denkt, durch $\frac{mm'}{rr}$ ausgedrückt. Das entspricht einer bewegenden Kraft $=\frac{mm'}{rr}$, die von der ersten Menge zur zweiten hin wirkt. Diese Formel gilt offenbar allgemein, wenn die Menge der südlichen Flüssigkeit negativ gerechnet wird, wie wir es von nun an stets voraussetzen. Ein negativer Wert der Kraft $\frac{mm'}{rr}$ bedeutet dann Anziehung anstelle von Abstoßung.

Befinden sich also in einem physischen Punkt zugleich gleiche Mengen nördlicher und südlicher Flüssigkeit, so geht davon überhaupt keine Wirkung aus. Sind die Mengen dagegen ungleich, so kommt nur der Überschuss der einen Flüssigkeit in Betracht. Diesen nennen wir \textit{freien} Magnetismus, positiv oder negativ.

\Article{2}
Zu diesen Grundannahmen müssen wir eine weitere hinzufügen, die durch die Erfahrung allenthalben bestätigt wird. Jeder Körper, in dem magnetische Flüssigkeiten vorhanden sind, enthält stets von beiden die gleiche Menge. Die Erfahrung lehrt sogar, dass diese Annahme auch für die einzelnen Teile eines solchen Körpers gilt, wie klein sie auch sein mögen, sofern unsere Sinne sie noch unterscheiden können. Wie wir am Ende des vorigen Artikels bemerkt haben, kann jedoch nur insofern eine Wirkung vorliegen, als eine gewisse Trennung der Flüssigkeiten stattfindet. Wir müssen also annehmen, dass die Abstände, über die sie sich trennen, so klein sind, dass sie sich unserer Messung entziehen.

Einen magnetisierbaren Körper muss man sich daher als Gefüge unzähliger Teilchen vorstellen, von denen jedes eine bestimmte Menge nördlicher und eine gleich große Menge südlicher magnetischer Flüssigkeit enthält. Die Flüssigkeiten sind entweder gleichmäßig vermischt, sodass der Magnetismus verborgen bleibt, oder sie haben sich mehr oder weniger voneinander getrennt, sodass der Magnetismus entwickelt ist. Diese Trennung kann jedoch niemals dazu führen, dass Flüssigkeit von \SourceLeaf{91}einem Teilchen in ein anderes übertritt. Es macht keinen Unterschied, ob man eine stärkere Trennung auf eine größere Menge freigewordener Flüssigkeiten oder auf einen größeren Abstand zwischen ihnen zurückführt. Bei der Größe der Trennung ist offenbar zugleich ihre Richtung zu berücksichtigen. Je nachdem, ob diese Richtungen in den verschiedenen Teilchen des Körpers übereinstimmen oder einander entgegenwirken, kann sich für Punkte außerhalb des Körpers eine größere oder kleinere Gesamtwirkung ergeben.

Wie der freie Magnetismus im Inneren des Körpers auch verteilt sein mag, nach einem allgemeineren Satz lässt sich gemäß einem bestimmten Gesetz stets eine andere, allein auf der Oberfläche des Körpers liegende Verteilung einsetzen, die hinsichtlich der nach außen wirkenden Kräfte der inneren Verteilung genau gleichwertig ist. Ein Element magnetischer Flüssigkeit erfährt dann an jedem Punkt außerhalb des Körpers von der tatsächlichen Verteilung des Magnetismus im Inneren genau dieselbe Anziehung oder Abstoßung wie von der gedachten Verteilung auf der Oberfläche. Dieselbe Fiktion lässt sich auf \textit{zwei} Körper ausdehnen, die durch den in ihnen entwickelten freien Magnetismus aufeinander wirken. Für jeden von beiden kann die gedachte Oberflächenverteilung an die Stelle der tatsächlichen inneren Verteilung treten. Erst auf diese Weise können wir der gewöhnlichen Redeweise, die etwa dem einen Ende einer Magnetnadel nur nördlichen und dem anderen südlichen Magnetismus zuschreibt, einen zutreffenden Sinn geben. Denn offenbar verträgt sich diese Ausdrucksweise nicht mit dem oben ausgesprochenen Grundsatz, den andere Erscheinungen zwingend verlangen. Es genüge, dies hier beiläufig zu bemerken. Auf den Satz selbst werden wir bei anderer Gelegenheit ausführlicher eingehen, da er für unser gegenwärtiges Vorhaben nicht erforderlich ist.

\Article{3}
Der \textit{magnetische Zustand} eines Körpers besteht in der Art, wie der freie Magnetismus in seinen einzelnen Teilchen verteilt ist. Was die Veränderlichkeit dieses Zustands betrifft, beobachten wir einen wesentlichen Unterschied zwischen den magnetisierbaren Körpern. Bei manchen, etwa bei weichem Eisen, ändert sich der Zustand schon durch die geringste Kraft sofort. Sobald die Kraft aufhört zu wirken, kehrt der frühere Zustand zurück. Bei anderen dagegen, insbesondere bei gehärtetem Stahl, muss die Kraft eine bestimmte Intensität erreicht haben, bevor sie eine merkliche Änderung des magnetischen Zustands hervorbringen kann. Hört die Kraft auf zu wirken, so bleibt der Körper entweder in dem erreichten Zustand oder kehrt wenigstens nicht ganz in den früheren zurück. In Körpern der ersten Art stellen sich die Moleküle der magnetischen Flüssigkeit also stets so ein, dass die Kräfte, die sie aufeinander ausüben und die aus äußeren Ursachen hervorgehen, im vollkommenen Gleichgewicht sind, oder sie weichen von einem solchen Gleichgewicht zumindest kaum merklich ab. In Körpern der zweiten Art dagegen kann der \SourceLeaf{92}magnetische Zustand auch ohne vollkommenes Gleichgewicht dieser Kräfte bestehen bleiben, sofern stärkere äußere Kräfte ferngehalten werden. Obwohl die Ursache dieser Erscheinung unbekannt ist, lässt sie sich so vorstellen, als setzten die wägbaren Teile eines Körpers der zweiten Art der Bewegung der mit ihnen verbundenen magnetischen Flüssigkeiten ein der Reibung ähnliches Hindernis entgegen. Bei weichem Eisen fehlt dieser Widerstand oder ist zumindest sehr gering.

In der theoretischen Untersuchung erfordern diese beiden Fälle eine ganz verschiedene Behandlung. In der vorliegenden Abhandlung wird aber nur von Körpern der zweiten Art die Rede sein. Bei den Versuchen, von denen wir sprechen werden, ist die Beständigkeit des magnetischen Zustands in jedem dafür verwendeten Körper eine Grundannahme. Man muss daher bei den Versuchen sorgfältig darauf achten, dass andere Körper, die diesen Zustand verändern könnten, nicht zu nahe kommen.

Es gibt jedoch eine Ursache für Änderungen, der auch Körper der zweiten Art unterliegen, nämlich die Wärme. Die Erfahrung zeigt, dass sich der magnetische Zustand eines Körpers mit seiner Temperatur ändert und dass zunehmende Wärme die Intensität des Magnetismus schwächt. Sofern der Körper nicht übermäßig erwärmt wurde, kehrt mit der früheren Temperatur auch der frühere magnetische Zustand zurück. Diese Abhängigkeit ist durch geeignete Versuche zu bestimmen. Wurden die zu einem Versuch gehörenden Vorgänge bei verschiedenen Temperaturen vorgenommen, müssen sie zunächst auf dieselbe Temperatur zurückgeführt werden.
